\subsection{Places and valuation rings}

PROPOSITION 1. Let K be a field, A a valuation ring of K and \(\kappa(A)\) the residue field of A. We extend the canonical mapping of A onto \(\kappa(A)\) to a mapping \(h_{A}: \tilde{K} \to (\kappa(A))^{\sim}\) by the equation \(h_{A}(x) = \infty\) if \(x \notin A\) . The mapping \(h_{A}\) thus defined is a place of K whose value field is \(\kappa(A)\) .

Clearly \(h_{\mathbf{A}}(1) = 1\) .

We show that \(h_{\mathbf{A}}\) is a morphism for addition. Let \(x, y\) be two elements of \(\tilde{\mathbf{K}}\) such that \(h_{\mathbf{A}}(x) + h_{\mathbf{A}}(y)\) is defined. One of the two elements \(x, y\) belongs then to \(\mathbf{A}\) and hence \(x + y\) is defined. If \(x \in \mathbf{A}\) and \(y \in \mathbf{A}\) , clearly

\[
h _ {\mathrm{A}} (x) + h _ {\mathrm{A}} (y) = h _ {\mathrm{A}} (x + y)
\]

holds. If \(x \in \mathbf{A}\) and \(y \notin \mathbf{A}\) , then \(x + y \notin \mathbf{A}\) and the two sides of the above formula equal \(\infty\) .

We show finally that \(h_{\mathbf{A}}\) is a morphism for multiplication. Let \(x \in \tilde{\mathbf{K}}, y \in \tilde{\mathbf{K}}\) be such that \(h_{\mathbf{A}}(x)h_{\mathbf{A}}(y)\) is defined. If \(x \in \mathbf{A}\) and \(y \in \mathbf{A}\) , clearly \(xy\) is defined and \(h_{\mathbf{A}}(x)h_{\mathbf{A}}(y) = h_{\mathbf{A}}(xy)\) . Suppose now that one of the elements \(x, y\) , for example \(y\) , does not belong to \(\mathbf{A}\) ; as \(h_{\mathbf{A}}(y) = \infty\) , \(h_{\mathbf{A}}(x) \neq 0\) , that is \(x \notin \mathfrak{m}(\mathbf{A})\) , whence \(x^{-1} \in \mathbf{A}\) ; it follows that \(xy\) is defined and that \(xy \notin \mathbf{A}\) , whence

\[
h _ {\mathrm{A}} (x y) = \infty = h _ {\mathrm{A}} (x) h _ {\mathrm{A}} (y).
\]

This proves Proposition 1.

If \(j\) is an isomorphism of \(\kappa(A)\) onto a subfield of a field \(L, j \circ h_{A}: \tilde{K} \to \tilde{L}\) is a place of \(K\) with values in \(L\) . The above process in fact provides all the places on \(K\) . To be precise:

PROPOSITION 2. Let K and L be two fields and f a place of K with values in L. Then there exists a valuation ring A of K and an isomorphism j of \(\kappa(\mathbf{A})\) onto a subfield of L such that \(f = j \circ h_{A}\) ; these conditions determine A and j uniquely. The ring A is the set of \(x \in K\) such that \(f(x) \neq \infty\) and \(\mathfrak{m}(\mathbf{A})\) is the set of \(x \in K\) such that \(f(x) = 0\) .

If \(f = j \circ h_{\mathbf{A}}\) , the condition \(f(x) \neq \infty\) (resp. \(f(x) = 0\) ) is equivalent to the condition \(h_{\mathbf{A}}(x) \neq \infty\) (resp. \(h_{\mathbf{A}}(x) = 0\) ) and hence to the condition \(x \in \mathbf{A}\) (resp. \(x \in \mathfrak{m}(\mathbf{A})\) ). Hence \(\mathbf{A}\) is determined uniquely and, as \(h_{\mathbf{A}}\) is surjective, \(j\) is also unique.

Now let \(f\) be any place of \(\mathbf{K}\) with values in \(\mathbf{L}\) ; let \(\mathbf{A}\) denote the set of \(x \in \mathbf{K}\) such that \(f(x) \neq \infty\) . If \(x \in \mathbf{A}\) and \(y \in \mathbf{A}\) , the compositions \(f(x) - f(y)\) and \(f(x)f(y)\) are defined and \(\neq \infty\) , which shows that \(x - y \in \mathbf{A}\) and \(xy \in \mathbf{A}\) ; hence \(\mathbf{A}\) is a subring of \(\mathbf{K}\) . If \(x \notin \mathbf{A}\) , then \(f(x) = \infty\) , hence \(f(x^{-1}) = 0\) and \(x^{-1}\) belongs to the kernel \(m\) of the homomorphism \(f'\) obtained by restricting \(f\) to \(\mathbf{A}\) . Conversely if \(y \in m\) , then \(y^{-1} \notin \mathbf{A}\) . This shows that \(\mathbf{A}\) is a valuation ring of \(\mathbf{K}\) and that \(m\) is its maximal ideal. Let \(j\) be the injective homomorphism from \(\kappa(\mathbf{A})\) to \(L\) derived from \(f'\) by passing to the quotient. Then \(f(x) = j(h_{\mathrm{A}}(x))\) for all \(x \in A\) and this equation remains true for \(x \notin A\) , the two sides being then equal to \(\infty\) .

The decomposition \(f = j \circ h_{A}\) is called the canonical decomposition of the place f.~A is called the ring of f, \(\mathfrak{m}(A)\) the ideal of f and \(\kappa(A)\) the residue field of f.~For two places \(f: \tilde{K} \to \tilde{L}\) and \(f': \tilde{K} \to \tilde{L}'\) to have the same ring, it is necessary and sufficient that there exist an isomorphism s of the value field of f onto that of \(f'\) such that \(f' = s \circ f\) ; then f and \(f'\) are called equivalent. It is seen that every result on valuation rings can be translated into a result on places and conversely; this is what we shall do in the following nos.

\textbf{Examples of places}\par

\begin{enumerate}
\def\labelenumi{(\arabic{enumi})}
\item
  Let \(\mathbf{K}\) be a field. The identity mapping on \(\mathbf{K}\) is a trivial place with ring \(\mathbf{K}\) and ideal (0).
\item
  Let \(k\) be a field. For all \(u \in k((\mathrm{T}))^{\sim}\) , let us write \(f(u) = \infty\) if \(u \notin k[[\mathrm{T}]]\) and define \(f(u)\) to be the constant term of \(u\) if \(u \in k[[\mathrm{T}]]\) . Then \(f\) is a place of \(k((\mathrm{T}))\) , with residue field \(k\) and ring \(k[[\mathrm{T}]]\) . For \(k[[\mathrm{T}]]\) is a valuation ring of \(k((\mathrm{T}))\) (§ 1, no. 4, Example 3) and the restriction of \(f\) to \(k[[\mathrm{T}]]\) is identified with the canonical homomorphism of \(k[[\mathrm{T}]]\) onto its residue field.
\item
  Let \(k\) be a field, \(a\) an element of \(k\) and \(A\) the set of \(u \in k(X)\) such that \(a\) is substitutable in \(u\) (Algebra, Chapter IV, § 3, no. 2). If \(p\) denotes the prime ideal \((X - a)\) of \(k[X]\) , then \(A = k[X]_p\) , so that \(A\) is a valuation ring of \(k(X)\) (§ 1, no. 4, Proposition 2). For all \(u \in k(X)^{\sim}\) , let us write \(f(u) = \infty\) if \(u \notin A\) and \(f(u) = u(a)\) if \(u \in A\) . Then \(f\) is a place of \(k(X)\) with residue field \(k\) and ring \(A\) ; for the restriction of \(f\) to \(A\) is a homomorphism of \(A\) onto \(k\) (Algebra, Chapter IV, § 3, Proposition 2) of kernel \(pA = m(A)\) . The element \(f(u) \in \tilde{k}\) is said to be obtained by putting \(X = a\) in \(u\) .
\end{enumerate}

* (4) Let S be a connected complex analytic variety of dimension 1 and K the field of meromorphic functions on S. For all \(z_0 \in \mathbf{S}\) , the mapping \(f \mapsto f(z_0)\) from K to \(\tilde{\mathbf{C}}\) is a place of K whose ring is the set of \(f \in \mathbf{K}\) which are holomorphic at \(z_0\) and whose ideal is the set of \(f \in \mathbf{K}\) which are zero at \(z_0\) . It is this example and other analogues which are the origin of the term ``place''.*

